\notechapter{N-20260416}{Line Integrals and Surface Integrals of the First and Second Kind}


\section{Line integrals with respect to arc length}

For a smooth curve $C:r(t)=(x(t),y(t))$, $a\le t\le b$, the first-kind line integral is
\[
  \int_C f(x,y)\,ds=\int_a^b f(x(t),y(t))\sqrt{x'(t)^2+y'(t)^2}\,dt.
\]
It integrates scalar density along length.

\section{Line integrals of vector fields}

For a vector field $F=(P,Q)$,
\[
  \int_C P\,dx+Q\,dy=\int_a^b [P(x(t),y(t))x'(t)+Q(x(t),y(t))y'(t)]\,dt.
\]
This is work or circulation.  Orientation matters.

\section{Surface integrals}

For a parametrized surface $r(u,v)$,
\[
  dS=\|r_u\times r_v\|\,du\,dv.
\]
The first-kind surface integral is
\[
  \iint_S f\,dS=\iint_D f(r(u,v))\|r_u\times r_v\|\,du\,dv.
\]
For a vector field $F$, the flux integral is
\[
  \iint_S F\cdot n\,dS=\iint_D F(r(u,v))\cdot(r_u\times r_v)\,du\,dv,
\]
with sign determined by orientation.

\section{Expanded synthesis and advanced context}

First-kind integrals integrate scalar density against geometric measure; second-kind integrals integrate vector fields through oriented curves or surfaces.  The difference is density versus flux.



\paragraph{Expanded synthesis.}
For this chapter, the elementary material should be read as a study of what remains invariant when coordinates change. The earlier sections (`Line integrals with respect to arc length`, `Line integrals of vector fields`, `Surface integrals`, `Higher viewpoint`, `Expanded distinction between first and second kind`) compute with arrays, bases, pivots, determinants, or eigenvectors; the deeper object is a linear map together with the spaces it acts on. A matrix is a representation of that map after bases have been chosen. This is why formulas such as
\[
  A' = P^{-1}AP,\qquad \widetilde A = T^{-1}AS
\]
are not cosmetic: they distinguish an intrinsic morphism from a coordinate description.

The structural hierarchy is as follows. Kernels record what a map forgets, images record what it reaches, cokernels record obstructions to solvability, and quotients record information after an equivalence has been imposed. Determinants act on the top exterior power and therefore measure oriented volume; traces are invariant under cyclic rearrangement and become characters in representation theory; adjoints depend on an inner product and lead to self-adjoint spectral theory.

A useful way to return to the computations is to ask, for every formula, which invariant it is measuring. Row reduction measures rank and kernel; diagonalization measures decomposition into invariant subspaces; Gram--Schmidt measures orthogonal projection; positive definiteness measures ellipsoid geometry; tensor notation measures how components transform. The elementary methods are therefore the finite-dimensional laboratory in which duality, quotienting, functoriality, and spectral decomposition can be seen clearly.


\section{Expanded distinction between first and second kind}

For line integrals, the first kind
\[
  \int_C f\,ds
\]
does not depend on orientation.  Reversing the parametrization leaves $ds$ unchanged.  The second kind
\[
  \int_C P\,dx+Q\,dy
\]
does depend on orientation, because $dx$ and $dy$ change sign.

For surface integrals, the same distinction holds.  The scalar surface integral $\iint_S f\,dS$ has no orientation.  The flux integral $\iint_SF\cdot n\,dS$ requires a choice of normal and changes sign when the normal is reversed.

\section{Parametric formulas as one template}

For curves, the geometric element is $|r'(t)|\,dt$.  For surfaces, it is $|r_u\times r_v|\,du\,dv$.  These are the same idea: the derivative of a parametrization tells how much the parameter domain is stretched into the actual geometric object.

\section{Elementary-to-advanced bridge: scalar integrals versus forms}

A first-kind line integral
\[
  \int_C f\,ds
\]
integrates a scalar density against arclength.  It depends on the metric because $ds$ is a length element.  A second-kind line integral
\[
  \int_C P\,dx+Q\,dy+R\,dz
\]
integrates a one-form.  It depends on orientation, not directly on arclength.

Similarly, a scalar surface integral
\[
  \iint_S f\,dS
\]
uses the metric area element, while a flux integral integrates a two-form or an oriented normal field.  This is why reversing orientation changes flux but not scalar area.

The elementary distinction between first and second kind is therefore the first encounter with the difference between metric measure and oriented differential forms.

\section{Elementary-to-advanced bridge: flux as a two-form}

For a vector field $F=(P,Q,R)$ in $\mathbb R^3$, the flux form is
\[
  \omega_F=P\,dy\wedge dz+Q\,dz\wedge dx+R\,dx\wedge dy.
\]
If $S$ is oriented, then
\[
  \iint_SF\cdot n\,dS=\iint_S \omega_F.
\]
The normal-vector formula is a Euclidean translation of this two-form integral.

This is more than notation.  Differential forms automatically encode orientation and pullback under parametrization.  If $r(u,v)$ parametrizes $S$, then
\[
  \iint_S\omega_F=\iint_D r^*\omega_F.
\]
The usual cross-product formula is the coordinate expression of this pullback.

\section{Elementary-to-advanced bridge: orientation bundles}

A surface may or may not admit a globally consistent normal direction.  Orientable surfaces, such as the sphere and torus, do.  The Möbius strip does not.  Flux integrals require an orientation, so they are globally well-defined only on oriented surfaces or after choosing local orientations carefully.

This explains why surface orientation is not a minor sign convention.  It is a topological property of the surface.

\section{Elementary-to-advanced bridge: manifolds with boundary}

A curve or surface used in integration is an example of a manifold.  A parametrized curve is a one-dimensional manifold; a surface is a two-dimensional manifold.  Boundaries appear naturally: a surface patch has boundary curves, and Stokes' theorem relates integration over the boundary to integration over the interior.

Thus the elementary parametrization formulas are local coordinate charts for manifolds.  The advanced language keeps the same computations but removes dependence on a single global parametrization.

% Second-pass deepening for waves 2-4: wave 4 part 2
\section{Second-pass deepening: parametrization invariance}

A well-defined geometric integral should not depend on the chosen parametrization.  For a first-kind line integral, reparametrizing $t=\phi(s)$ changes $|r'(t)|dt$ into $|(r\circ\phi)'(s)|ds$, so the value is unchanged for orientation-preserving reparametrizations.

For a one-form integral, orientation matters.  Reversing the parametrization changes sign.  This is exactly what should happen for work or circulation: traversing a path backward reverses the accumulated oriented effect.

\section{Second-pass deepening: surface normals as choices of orientation}

For a parametrized surface $r(u,v)$, the vector $r_u\times r_v$ depends on the order of parameters.  Swapping $u$ and $v$ reverses orientation.  The flux integral is therefore an integral over an oriented surface, not merely over a set of points.

The elementary cross-product normal formula is the coordinate expression of choosing an orientation on a two-dimensional manifold.
